% =============================================================================
% ESL FRAMEWORK - CHAPTER 4: LINGUISTIC MARKERS
% =============================================================================
% Authors: Gerhard Fehr & Andrea Fehr
% FehrAdvice & Partners AG, Zürich
% Version: Extended Monograph
% =============================================================================

\documentclass[11pt,a4paper]{article}

% --- Packages ---
\usepackage[utf8]{inputenc}
\usepackage[T1]{fontenc}
\usepackage[english]{babel}
\usepackage{amsmath,amssymb,amsthm}
\usepackage{mathtools}
\usepackage{booktabs}
\usepackage{longtable}
\usepackage{array}
\usepackage{hyperref}
\usepackage{geometry}
\usepackage{xcolor}
\usepackage{tcolorbox}
\usepackage{tikz}
\usepackage{pgfplots}
\usepackage{natbib}
\usepackage{enumitem}
\usepackage{fancyhdr}
\usepackage{titlesec}
\usepackage{multirow}
\usepackage{colortbl}
\pgfplotsset{compat=1.17}
\usetikzlibrary{shapes,arrows.meta,positioning,calc,patterns}

\geometry{margin=2.5cm}

% --- Colors ---
\definecolor{eslblue}{RGB}{0,82,147}
\definecolor{eslgreen}{RGB}{0,128,0}
\definecolor{eslorange}{RGB}{230,159,0}
\definecolor{eslred}{RGB}{204,51,17}
\definecolor{eslgray}{RGB}{128,128,128}
\definecolor{eslpurple}{RGB}{148,0,211}
\definecolor{strongcolor}{RGB}{180,0,0}
\definecolor{moderatecolor}{RGB}{200,150,0}
\definecolor{weakcolor}{RGB}{0,120,180}

% --- Boxes ---
\newtcolorbox{eslbox}[1][]{
  colback=eslblue!5!white,
  colframe=eslblue,
  title={ESL Assessment},
  fonttitle=\bfseries,
  #1
}

\newtcolorbox{keyinsight}[1][]{
  colback=eslgreen!5!white,
  colframe=eslgreen!80!black,
  title={Key Insight},
  fonttitle=\bfseries,
  #1
}

\newtcolorbox{warningbox}[1][]{
  colback=eslorange!5!white,
  colframe=eslorange!80!black,
  title={Validation Status},
  fonttitle=\bfseries,
  #1
}

\newtcolorbox{examplebox}[1][]{
  colback=eslgray!10!white,
  colframe=eslgray!80!black,
  title={Example},
  fonttitle=\bfseries,
  #1
}

\newtcolorbox{definitionbox}[1][]{
  colback=white,
  colframe=eslblue!80!black,
  title={Definition},
  fonttitle=\bfseries,
  #1
}

\newtcolorbox{lexiconbox}[1][]{
  colback=eslpurple!5!white,
  colframe=eslpurple!80!black,
  title={Marker Lexicon},
  fonttitle=\bfseries,
  #1
}

\newtcolorbox{calculationbox}[1][]{
  colback=eslblue!5!white,
  colframe=eslblue!80!black,
  title={Calculation},
  fonttitle=\bfseries,
  #1
}

% --- Theorems ---
\newtheorem{definition}{Definition}[section]
\newtheorem{theorem}{Theorem}[section]
\newtheorem{principle}{Principle}[section]
\newtheorem{proposition}[theorem]{Proposition}
\newtheorem{rule_def}{Rule}[section]

\theoremstyle{remark}
\newtheorem{remark}{Remark}[section]

% --- Header/Footer ---
\pagestyle{fancy}
\fancyhf{}
\fancyhead[L]{\small ESL Framework -- Extended Monograph}
\fancyhead[R]{\small Chapter 4: Linguistic Markers}
\fancyfoot[C]{\thepage}

% =============================================================================
% TITLE
% =============================================================================

\title{%
    \Large \textbf{Empirical Stability of Language (ESL)} \\[0.5em]
    \large A Universal Quality Assurance Framework for \\
    Scientific Claim Validation \\[1em]
    \LARGE \textbf{Chapter 4: Linguistic Markers} \\[0.3em]
    \large B-Value Operationalization and Claim Strength Assessment
}

\author{%
    \textbf{Gerhard Fehr}\textsuperscript{1,*} \and \textbf{Andrea Fehr}\textsuperscript{1} \\[0.5em]
    \textsuperscript{1}FehrAdvice \& Partners AG, Binzmühlestrasse 170A, 8050 Zürich, Switzerland \\[0.3em]
    \textsuperscript{*}Corresponding author: andrea.fehr@fehradvice.com
}

\date{December 26, 2025}

\begin{document}

\maketitle

% =============================================================================
% ABSTRACT
% =============================================================================

\begin{abstract}
\noindent
This chapter develops the B-component of the K-metric: the systematic identification and quantification of linguistic markers that signal claim strength. We present a comprehensive marker lexicon organized into five categories: epistemic verbs (proves, demonstrates, suggests), modal auxiliaries (must, may, might), quantifiers (all, most, some), hedging expressions (preliminary, in this sample), and intensifiers/downtoners. For each category, we provide proposed B-values based on ordinal linguistic analysis, with explicit acknowledgment of validation status. The chapter introduces the Composite B-Calculation algorithm for statements containing multiple markers, addresses context-sensitivity and discipline-specific variation, and outlines the Monte-Carlo validation approach for empirical B-value calibration. We conclude with practical guidelines for B-assessment in scientific texts and LLM-generated content.

\vspace{0.5em}
\noindent\textbf{Keywords:} Linguistic markers, epistemic modality, hedging, claim strength, B-values, corpus linguistics
\end{abstract}

\tableofcontents
\newpage

% =============================================================================
% 4.1 THEORETICAL FOUNDATIONS
% =============================================================================

\section{Theoretical Foundations}
\label{sec:foundations}

\subsection{Epistemic Modality in Language}

Human languages have evolved sophisticated systems for expressing degrees of certainty. Linguists refer to this as \textbf{epistemic modality}---the grammatical and lexical resources for indicating a speaker's confidence in a proposition.

\begin{definition}[Epistemic Modality]
\label{def:epistemic}
Epistemic modality is the linguistic expression of the speaker's degree of commitment to the truth of a proposition. It encompasses:
\begin{itemize}
    \item \textbf{Modal verbs}: must, may, might, could, should
    \item \textbf{Epistemic adverbs}: certainly, probably, possibly, perhaps
    \item \textbf{Evidential markers}: apparently, reportedly, allegedly
    \item \textbf{Hedging expressions}: somewhat, to some extent, in certain cases
\end{itemize}
\end{definition}

\subsection{The Linguistics Literature}

ESL's B-value framework draws on established corpus linguistics research:

\begin{itemize}
    \item \textbf{Coates (1983)}: Systematic analysis of English modal verbs, establishing that modals form a cline from strong (must) to weak (might).

    \item \textbf{Hyland (1998)}: Comprehensive study of hedging in academic writing, documenting pervasive uncertainty marking in scientific discourse.

    \item \textbf{Salager-Meyer (1994)}: Analysis of hedging in medical research articles, showing discipline-specific patterns.

    \item \textbf{Biber et al. (1999)}: Longman Grammar of Spoken and Written English, providing corpus-based frequency data for epistemic markers.
\end{itemize}

\begin{keyinsight}[title={The Linguistic Foundation}]
The \textit{existence} of epistemic markers in language is well-established (Claim Type T1). The \textit{specific B-values} assigned to markers are proposed calibrations requiring empirical validation (Claim Type T4--T5).
\end{keyinsight}

\subsection{From Ordinal Ranking to Cardinal Values}

Linguistics establishes ordinal rankings (``must'' is stronger than ``may''). ESL requires cardinal values ($B_{\text{must}} = 0.95$, $B_{\text{may}} = 0.40$).

\begin{principle}[Ordinal-to-Cardinal Mapping]
\label{princ:mapping}
B-value assignment follows three principles:
\begin{enumerate}
    \item \textbf{Preserve ordinal ranking}: If $A$ is linguistically stronger than $B$, then $B_A > B_B$.
    \item \textbf{Span the interval}: Strongest markers approach 1.0; weakest approach 0.0.
    \item \textbf{Maintain discriminability}: Adjacent markers have distinguishable B-values ($|\Delta B| \geq 0.05$).
\end{enumerate}
\end{principle}

\begin{warningbox}
The specific numerical B-values in this chapter are \textbf{proposed estimates} based on:
\begin{itemize}
    \item Ordinal linguistic rankings (empirically grounded)
    \item Expert judgment (author intuition)
    \item Internal consistency requirements
\end{itemize}

\textbf{Validation status}: T4 (Expert Proposal) to T5 (Author Intuition). Corpus validation studies are needed to establish empirical B-values. See Section \ref{sec:validation} for the Monte-Carlo validation approach.
\end{warningbox}

% =============================================================================
% 4.2 MARKER CATEGORY 1: EPISTEMIC VERBS
% =============================================================================

\section{Epistemic Verbs}
\label{sec:epistemic-verbs}

Epistemic verbs are the primary carriers of claim strength in scientific writing. They signal what the evidence ``does''---proves, demonstrates, suggests, or merely hints.

\subsection{The Epistemic Verb Hierarchy}

\begin{lexiconbox}[title={Epistemic Verb B-Values}]
\begin{center}
\begin{tabular}{llcc}
\toprule
\textbf{Category} & \textbf{Verb/Phrase} & \textbf{B-Value} & \textbf{95\% Range} \\
\midrule
\multirow{4}{*}{\textcolor{strongcolor}{\textbf{Definitive}}}
    & proves, prove & 0.95 & [0.92, 0.98] \\
    & demonstrates conclusively & 0.94 & [0.91, 0.97] \\
    & establishes & 0.92 & [0.88, 0.95] \\
    & confirms & 0.90 & [0.86, 0.93] \\
\midrule
\multirow{4}{*}{\textcolor{strongcolor}{\textbf{Strong}}}
    & demonstrates & 0.88 & [0.84, 0.91] \\
    & shows & 0.85 & [0.80, 0.88] \\
    & reveals & 0.83 & [0.78, 0.87] \\
    & finds & 0.80 & [0.75, 0.84] \\
\midrule
\multirow{4}{*}{\textcolor{moderatecolor}{\textbf{Moderate}}}
    & indicates & 0.72 & [0.68, 0.76] \\
    & supports & 0.70 & [0.65, 0.74] \\
    & suggests & 0.62 & [0.58, 0.66] \\
    & points to & 0.60 & [0.55, 0.64] \\
\midrule
\multirow{4}{*}{\textcolor{weakcolor}{\textbf{Tentative}}}
    & is consistent with & 0.55 & [0.50, 0.60] \\
    & is associated with & 0.52 & [0.47, 0.57] \\
    & hints at & 0.45 & [0.40, 0.50] \\
    & raises the possibility & 0.42 & [0.37, 0.47] \\
\midrule
\multirow{3}{*}{\textcolor{weakcolor}{\textbf{Minimal}}}
    & does not rule out & 0.35 & [0.30, 0.40] \\
    & cannot exclude & 0.32 & [0.27, 0.37] \\
    & may be related to & 0.30 & [0.25, 0.35] \\
\bottomrule
\end{tabular}
\end{center}
\end{lexiconbox}

\subsection{Usage Patterns}

\begin{examplebox}[title={Epistemic Verb Comparison}]
The same finding can be expressed with different claim strengths:

\begin{enumerate}
    \item ``Our results \textbf{prove} that X causes Y.'' ($B = 0.95$)
    \item ``Our results \textbf{demonstrate} that X causes Y.'' ($B = 0.88$)
    \item ``Our results \textbf{suggest} that X causes Y.'' ($B = 0.62$)
    \item ``Our results \textbf{are consistent with} X causing Y.'' ($B = 0.55$)
    \item ``Our results \textbf{do not rule out} that X causes Y.'' ($B = 0.35$)
\end{enumerate}

Each formulation conveys a different degree of epistemic commitment. Calibrated writing matches verb choice to evidence strength.
\end{examplebox}

\subsection{Context Sensitivity}

Epistemic verb strength varies by context:

\begin{table}[h]
\centering
\caption{Context-Dependent B-Value Adjustments}
\label{tab:context-verbs}
\begin{tabular}{lcc}
\toprule
\textbf{Context} & \textbf{Adjustment} & \textbf{Rationale} \\
\midrule
Abstract/Conclusion & $+0.00$ & Base interpretation \\
Methods section & $-0.05$ & Procedural, not evidential \\
Literature review & $-0.03$ & Reporting others' claims \\
Speculation section & $-0.08$ & Explicitly tentative \\
Title & $+0.05$ & Headline effect \\
\bottomrule
\end{tabular}
\end{table}

% =============================================================================
% 4.3 MARKER CATEGORY 2: MODAL AUXILIARIES
% =============================================================================

\section{Modal Auxiliaries}
\label{sec:modals}

Modal auxiliary verbs (must, will, should, may, might, could) are core grammatical resources for epistemic modality in English.

\subsection{The Modal Hierarchy}

\begin{lexiconbox}[title={Modal Auxiliary B-Values}]
\begin{center}
\begin{tabular}{llcc}
\toprule
\textbf{Strength} & \textbf{Modal} & \textbf{B-Value} & \textbf{Example} \\
\midrule
\multirow{2}{*}{\textcolor{strongcolor}{\textbf{Necessity}}}
    & must (epistemic) & 0.92 & ``This must be the cause'' \\
    & has to & 0.90 & ``It has to follow that...'' \\
\midrule
\multirow{2}{*}{\textcolor{strongcolor}{\textbf{Confidence}}}
    & will (prediction) & 0.85 & ``This will replicate'' \\
    & would (conditional) & 0.78 & ``This would suggest...'' \\
\midrule
\multirow{2}{*}{\textcolor{moderatecolor}{\textbf{Probability}}}
    & should & 0.72 & ``Results should generalize'' \\
    & ought to & 0.70 & ``We ought to observe...'' \\
\midrule
\multirow{3}{*}{\textcolor{weakcolor}{\textbf{Possibility}}}
    & may & 0.45 & ``This may indicate...'' \\
    & might & 0.40 & ``It might be that...'' \\
    & could & 0.38 & ``This could explain...'' \\
\midrule
\textcolor{weakcolor}{\textbf{Remote}}
    & can (theoretical) & 0.35 & ``One can imagine...'' \\
\bottomrule
\end{tabular}
\end{center}
\end{lexiconbox}

\subsection{Epistemic vs. Deontic Readings}

Modal verbs have multiple meanings. ESL focuses on epistemic (knowledge-related) readings:

\begin{definition}[Epistemic vs. Deontic Modality]
\label{def:epistemic-deontic}
\begin{itemize}
    \item \textbf{Epistemic}: Relates to knowledge or belief. \\
    ``You \textit{must} be tired'' = I conclude you are tired.

    \item \textbf{Deontic}: Relates to obligation or permission. \\
    ``You \textit{must} submit by Friday'' = You are required to submit.
\end{itemize}
\end{definition}

\begin{rule_def}[Modal Disambiguation]
\label{rule:modal}
In scientific writing, prioritize epistemic readings unless context clearly indicates deontic meaning (e.g., methodological requirements).
\end{rule_def}

% =============================================================================
% 4.4 MARKER CATEGORY 3: QUANTIFIERS
% =============================================================================

\section{Quantifiers}
\label{sec:quantifiers}

Quantifiers modulate claim scope---the range of cases to which a claim applies.

\subsection{The Quantifier Hierarchy}

\begin{lexiconbox}[title={Quantifier B-Values}]
\begin{center}
\begin{tabular}{llcc}
\toprule
\textbf{Scope} & \textbf{Quantifier} & \textbf{B-Value} & \textbf{Claim Scope} \\
\midrule
\multirow{4}{*}{\textcolor{strongcolor}{\textbf{Universal}}}
    & all, every & 0.98 & 100\% of cases \\
    & always, invariably & 0.96 & No exceptions \\
    & never, none & 0.96 & Universal negation \\
    & without exception & 0.98 & Explicit universality \\
\midrule
\multirow{3}{*}{\textcolor{strongcolor}{\textbf{Near-Universal}}}
    & almost all & 0.90 & $>$95\% of cases \\
    & nearly always & 0.88 & Few exceptions \\
    & the vast majority & 0.86 & $>$90\% of cases \\
\midrule
\multirow{4}{*}{\textcolor{moderatecolor}{\textbf{Majority}}}
    & most & 0.78 & $>$50\% of cases \\
    & typically, usually & 0.75 & Normal/expected \\
    & generally & 0.72 & With exceptions \\
    & often & 0.68 & Frequent \\
\midrule
\multirow{4}{*}{\textcolor{moderatecolor}{\textbf{Common}}}
    & many & 0.62 & Substantial portion \\
    & frequently & 0.60 & Regular occurrence \\
    & commonly & 0.58 & Not unusual \\
    & tends to & 0.55 & Probabilistic \\
\midrule
\multirow{4}{*}{\textcolor{weakcolor}{\textbf{Partial}}}
    & some & 0.45 & Unspecified portion \\
    & sometimes & 0.42 & Occasional \\
    & in some cases & 0.40 & Restricted scope \\
    & occasionally & 0.38 & Infrequent \\
\midrule
\multirow{3}{*}{\textcolor{weakcolor}{\textbf{Rare}}}
    & few & 0.28 & Small minority \\
    & rarely, seldom & 0.25 & Uncommon \\
    & in rare cases & 0.22 & Exceptional \\
\bottomrule
\end{tabular}
\end{center}
\end{lexiconbox}

\subsection{Quantifier Scope Interactions}

\begin{examplebox}[title={Quantifier Impact on Claim Strength}]
Consider variations of a finding:

\begin{enumerate}
    \item ``Loss aversion affects \textbf{all} economic decisions.'' ($B = 0.98$)
    \item ``Loss aversion affects \textbf{most} economic decisions.'' ($B = 0.78$)
    \item ``Loss aversion affects \textbf{many} economic decisions.'' ($B = 0.62$)
    \item ``Loss aversion affects \textbf{some} economic decisions.'' ($B = 0.45$)
    \item ``Loss aversion \textbf{rarely} affects economic decisions.'' ($B = 0.25$)
\end{enumerate}

The quantifier alone shifts B by up to 0.73 points.
\end{examplebox}

% =============================================================================
% 4.5 MARKER CATEGORY 4: HEDGING EXPRESSIONS
% =============================================================================

\section{Hedging Expressions}
\label{sec:hedging}

Hedging expressions explicitly signal limitations, tentativeness, or reduced scope. Unlike epistemic verbs, hedges typically function as \textit{modifiers} that reduce the B-value of the primary claim.

\subsection{Hedge Categories}

\begin{lexiconbox}[title={Hedging Expression Modifiers ($\Delta B$)}]
\begin{center}
\begin{tabular}{llc}
\toprule
\textbf{Category} & \textbf{Expression} & \textbf{$\Delta B$} \\
\midrule
\multirow{4}{*}{\textbf{Tentativeness}}
    & preliminary, initial & $-0.15$ \\
    & exploratory & $-0.18$ \\
    & tentative & $-0.20$ \\
    & speculative & $-0.25$ \\
\midrule
\multirow{4}{*}{\textbf{Scope Limitation}}
    & in this sample/study & $-0.10$ \\
    & in this population & $-0.10$ \\
    & under these conditions & $-0.12$ \\
    & in this context & $-0.08$ \\
\midrule
\multirow{3}{*}{\textbf{Methodological}}
    & with the caveat that... & $-0.10$ \\
    & subject to limitations & $-0.12$ \\
    & with reservations & $-0.15$ \\
\midrule
\multirow{4}{*}{\textbf{Future-Oriented}}
    & further research is needed & $-0.18$ \\
    & requires replication & $-0.15$ \\
    & awaits confirmation & $-0.20$ \\
    & remains to be seen & $-0.22$ \\
\midrule
\multirow{3}{*}{\textbf{Approximators}}
    & approximately, roughly & $-0.05$ \\
    & somewhat, to some extent & $-0.08$ \\
    & more or less & $-0.10$ \\
\bottomrule
\end{tabular}
\end{center}
\end{lexiconbox}

\subsection{Cumulative Hedging}

Multiple hedges accumulate:

\begin{examplebox}[title={Cumulative Hedging Effect}]
\textbf{Statement}: ``\textit{Preliminary} results \textit{in this sample} \textit{tentatively} suggest that, \textit{subject to limitations}, X may affect Y, though \textit{further research is needed}.''

\textbf{Calculation}:
\begin{align}
B_{\text{base}} &= 0.62 \quad \text{(``suggest'')} \\
\Delta_{\text{preliminary}} &= -0.15 \\
\Delta_{\text{in this sample}} &= -0.10 \\
\Delta_{\text{tentatively}} &= -0.20 \\
\Delta_{\text{subject to limitations}} &= -0.12 \\
\Delta_{\text{further research}} &= -0.18 \\
\hline
B_{\text{composite}} &= 0.62 - 0.15 - 0.10 - 0.20 - 0.12 - 0.18 \\
&= 0.62 - 0.75 = -0.13 \to 0.00 \text{ (floor)}
\end{align}

Excessive hedging produces near-zero claim strength.
\end{examplebox}

% =============================================================================
% 4.6 MARKER CATEGORY 5: INTENSIFIERS AND DOWNTONERS
% =============================================================================

\section{Intensifiers and Downtoners}
\label{sec:intensifiers}

Adverbial modifiers adjust claim strength up (intensifiers) or down (downtoners).

\subsection{Intensifiers}

\begin{table}[h]
\centering
\caption{Intensifier B-Value Modifiers}
\label{tab:intensifiers}
\begin{tabular}{lcc}
\toprule
\textbf{Intensifier} & \textbf{$\Delta B$} & \textbf{Example} \\
\midrule
clearly, obviously & $+0.08$ & ``clearly demonstrates'' \\
strongly, powerfully & $+0.10$ & ``strongly suggests'' \\
definitively, conclusively & $+0.12$ & ``definitively establishes'' \\
unequivocally, unambiguously & $+0.15$ & ``unequivocally confirms'' \\
beyond doubt & $+0.18$ & ``proves beyond doubt'' \\
\bottomrule
\end{tabular}
\end{table}

\subsection{Downtoners}

\begin{table}[h]
\centering
\caption{Downtoner B-Value Modifiers}
\label{tab:downtoners}
\begin{tabular}{lcc}
\toprule
\textbf{Downtoner} & \textbf{$\Delta B$} & \textbf{Example} \\
\midrule
slightly, marginally & $-0.08$ & ``slightly suggests'' \\
partially, partly & $-0.10$ & ``partially supports'' \\
weakly & $-0.12$ & ``weakly indicates'' \\
barely, hardly & $-0.15$ & ``barely supports'' \\
\bottomrule
\end{tabular}
\end{table}

% =============================================================================
% 4.7 COMPOSITE B-CALCULATION
% =============================================================================

\section{Composite B-Calculation}
\label{sec:composite}

Real scientific statements contain multiple markers. ESL provides an algorithm for computing composite B-values.

\subsection{The Composite B Algorithm}

\begin{definition}[Composite B-Calculation]
\label{def:composite-b}
For a statement with primary marker $m_0$ and modifiers $m_1, \ldots, m_n$:

\begin{equation}
B_{\text{composite}} = \max\left(0, \min\left(1, B_{m_0} + \sum_{i=1}^{n} \Delta B_{m_i}\right)\right)
\end{equation}

where:
\begin{itemize}
    \item $B_{m_0}$ is the base B-value of the primary epistemic verb
    \item $\Delta B_{m_i}$ is the modifier effect of marker $i$
    \item Results are bounded to $[0, 1]$
\end{itemize}
\end{definition}

\subsection{Marker Identification Priority}

\begin{rule_def}[Marker Priority]
\label{rule:priority}
When identifying markers in a statement:
\begin{enumerate}
    \item \textbf{Primary verb}: Identify the main epistemic verb (proves, suggests, etc.)
    \item \textbf{Modal auxiliary}: Check for modals modifying the verb
    \item \textbf{Quantifiers}: Identify scope-modifying quantifiers
    \item \textbf{Hedges}: Count all hedging expressions
    \item \textbf{Intensifiers/Downtoners}: Apply adverbial modifiers
\end{enumerate}
\end{rule_def}

\subsection{Worked Examples}

\begin{calculationbox}[title={Example 1: Strong Claim}]
\textbf{Statement}: ``Our experiments \textbf{conclusively demonstrate} that all subjects exhibit loss aversion.''

\textbf{Markers}:
\begin{itemize}
    \item Primary verb: ``demonstrate'' ($B = 0.88$)
    \item Intensifier: ``conclusively'' ($\Delta B = +0.12$)
    \item Quantifier: ``all'' ($B_{\text{quant}} = 0.98$, but scope marker)
\end{itemize}

\textbf{Calculation}:
\begin{equation}
B = 0.88 + 0.12 = 1.00 \to 0.98 \text{ (adjusted for universal claim)}
\end{equation}
\end{calculationbox}

\begin{calculationbox}[title={Example 2: Moderate Claim}]
\textbf{Statement}: ``Results \textbf{suggest} that loss aversion \textbf{may} affect \textbf{many} decisions.''

\textbf{Markers}:
\begin{itemize}
    \item Primary verb: ``suggest'' ($B = 0.62$)
    \item Modal: ``may'' ($\Delta B = -0.17$, i.e., moves toward 0.45)
    \item Quantifier: ``many'' (scope effect, not additive)
\end{itemize}

\textbf{Calculation} (using weighted average for modal interaction):
\begin{equation}
B = 0.62 \times 0.6 + 0.45 \times 0.4 = 0.372 + 0.180 = 0.55
\end{equation}

Alternative (simpler additive):
\begin{equation}
B = 0.62 - 0.10 = 0.52 \text{ (modal penalty)}
\end{equation}
\end{calculationbox}

\begin{calculationbox}[title={Example 3: Heavily Hedged Claim}]
\textbf{Statement}: ``\textbf{Preliminary} analysis of \textbf{this sample} \textbf{tentatively suggests} that X \textbf{might} be associated with Y, \textbf{subject to limitations} and \textbf{pending further research}.''

\textbf{Markers}:
\begin{itemize}
    \item Primary verb: ``suggests'' ($B = 0.62$)
    \item Hedge: ``preliminary'' ($\Delta B = -0.15$)
    \item Hedge: ``this sample'' ($\Delta B = -0.10$)
    \item Hedge: ``tentatively'' ($\Delta B = -0.20$)
    \item Modal: ``might'' ($\Delta B = -0.10$)
    \item Hedge: ``subject to limitations'' ($\Delta B = -0.12$)
    \item Hedge: ``pending further research'' ($\Delta B = -0.18$)
\end{itemize}

\textbf{Calculation}:
\begin{align}
B &= 0.62 - 0.15 - 0.10 - 0.20 - 0.10 - 0.12 - 0.18 \\
&= 0.62 - 0.85 = -0.23 \to 0.00 \text{ (floor)}
\end{align}

\textbf{Interpretation}: The claim is so heavily hedged that it conveys essentially no epistemic commitment.
\end{calculationbox}

% =============================================================================
% 4.8 DISCIPLINE-SPECIFIC VARIATION
% =============================================================================

\section{Discipline-Specific Variation}
\label{sec:discipline-variation}

B-values are not universal constants. The same word may signal different claim strengths across disciplines.

\subsection{Cross-Disciplinary Calibration}

\begin{table}[h]
\centering
\caption{Discipline-Specific B-Value Variation}
\label{tab:discipline-b}
\begin{tabular}{lcccc}
\toprule
\textbf{Marker} & \textbf{Physics} & \textbf{Medicine} & \textbf{Psychology} & \textbf{Law} \\
\midrule
``proves'' & 0.98 & 0.90 & 0.88 & 0.95 \\
``demonstrates'' & 0.92 & 0.85 & 0.82 & 0.88 \\
``suggests'' & 0.70 & 0.60 & 0.58 & 0.65 \\
``is consistent with'' & 0.65 & 0.52 & 0.50 & 0.58 \\
\bottomrule
\end{tabular}
\end{table}

\subsection{Disciplinary Defaults}

\begin{principle}[Discipline Default Rule]
\label{princ:discipline-default}
In absence of explicit discipline specification:
\begin{enumerate}
    \item Use ``General Science'' B-values (this chapter's defaults)
    \item Apply discipline-specific adjustments when known
    \item Document any adjustments in ESL assessment
\end{enumerate}
\end{principle}

% =============================================================================
% 4.9 MONTE-CARLO VALIDATION APPROACH
% =============================================================================

\section{Monte-Carlo Validation Approach}
\label{sec:validation}

B-values in this chapter are expert proposals. ESL 2.1 introduces Monte-Carlo validation for empirical calibration.

\subsection{The Validation Problem}

\begin{definition}[B-Value Validation]
\label{def:b-validation}
A B-value is \textbf{empirically validated} when:
\begin{enumerate}
    \item Multiple independent raters assign similar values (inter-rater reliability)
    \item Values predict outcomes (predictive validity)
    \item Values are stable across samples (test-retest reliability)
\end{enumerate}
\end{definition}

\subsection{Monte-Carlo LLM Sampling}

ESL 2.1 proposes using LLM sampling for B-value estimation:

\begin{definition}[Monte-Carlo B-Estimation]
\label{def:mc-b}
For a linguistic marker $m$:
\begin{equation}
\hat{B}_m = \frac{1}{N} \sum_{i=1}^{N} B_m^{(i)}
\end{equation}
where $B_m^{(i)}$ is the $i$-th LLM estimate under varied conditions (temperature, prompt variant).
\end{definition}

\begin{definition}[Consistency Coefficient]
\label{def:kappa-mc}
\begin{equation}
\kappa_{\text{MC}} = 1 - \frac{\sigma_m}{\sigma_{\max}}
\end{equation}
where $\sigma_{\max} = 0.29$ (standard deviation of uniform distribution on $[0,1]$).

High $\kappa_{\text{MC}}$ indicates consistent LLM estimates; low $\kappa_{\text{MC}}$ indicates marker ambiguity.
\end{definition}

\subsection{Validation Status Summary}

\begin{table}[h]
\centering
\caption{B-Value Validation Status by Category}
\label{tab:validation-status}
\begin{tabular}{lccc}
\toprule
\textbf{Category} & \textbf{Ordinal Ranking} & \textbf{Cardinal Values} & \textbf{Claim Type} \\
\midrule
Epistemic verbs & T1 (Empirical) & T4--T5 (Proposal) & Partial \\
Modal auxiliaries & T1 (Empirical) & T4--T5 (Proposal) & Partial \\
Quantifiers & T2 (Established) & T4 (Proposal) & Partial \\
Hedging expressions & T2 (Established) & T5 (Intuition) & Weak \\
Intensifiers & T3 (Logical) & T5 (Intuition) & Weak \\
\bottomrule
\end{tabular}
\end{table}

% =============================================================================
% 4.10 PRACTICAL GUIDELINES
% =============================================================================

\section{Practical Guidelines for B-Assessment}
\label{sec:guidelines}

\subsection{Step-by-Step B-Assessment}

\begin{enumerate}
    \item \textbf{Identify the core claim}: What proposition is being asserted?

    \item \textbf{Find the primary epistemic verb}: What verb links evidence to claim? (proves, suggests, indicates, etc.)

    \item \textbf{Check for modal auxiliaries}: Is the claim modified by must, may, might, could?

    \item \textbf{Identify quantifiers}: What scope does the claim have? (all, most, some, etc.)

    \item \textbf{Count hedging expressions}: How many qualifications, limitations, or caveats are present?

    \item \textbf{Apply intensifiers/downtoners}: Are there adverbial modifiers?

    \item \textbf{Calculate composite B}: Use the additive formula with floor/ceiling bounds.

    \item \textbf{Document your assessment}: Record markers identified and calculation steps.
\end{enumerate}

\subsection{Common Pitfalls}

\begin{table}[h]
\centering
\caption{Common B-Assessment Errors}
\label{tab:pitfalls}
\begin{tabular}{lp{8cm}}
\toprule
\textbf{Pitfall} & \textbf{Solution} \\
\midrule
Ignoring hedges & Count all qualifications systematically \\
Double-counting & Each marker applies once \\
Context neglect & Adjust for section type (abstract vs. methods) \\
Discipline blindness & Apply discipline-specific calibration \\
Missing implicit hedges & ``Initial,'' ``pilot,'' ``exploratory'' are hedges \\
\bottomrule
\end{tabular}
\end{table}

\subsection{Quick Reference Card}

\begin{table}[h]
\centering
\caption{B-Value Quick Reference}
\label{tab:quick-ref}
\begin{tabular}{lcc}
\toprule
\textbf{Claim Strength} & \textbf{B Range} & \textbf{Typical Markers} \\
\midrule
Definitive & 0.90--1.00 & proves, establishes, confirms \\
Strong & 0.75--0.89 & demonstrates, shows, reveals \\
Moderate & 0.55--0.74 & indicates, suggests, supports \\
Tentative & 0.35--0.54 & is consistent with, may indicate \\
Minimal & 0.15--0.34 & does not rule out, hints at \\
Speculative & 0.00--0.14 & (heavily hedged claims) \\
\bottomrule
\end{tabular}
\end{table}

% =============================================================================
% 4.11 CHAPTER SUMMARY
% =============================================================================

\section{Chapter Summary}
\label{sec:summary}

This chapter has developed the B-component of the K-metric:

\begin{enumerate}
    \item \textbf{Theoretical Foundation}: Epistemic modality in language is well-established; specific B-values are proposed calibrations.

    \item \textbf{Five Marker Categories}:
    \begin{itemize}
        \item Epistemic verbs (proves → hints at): $B \in [0.30, 0.95]$
        \item Modal auxiliaries (must → might): $B \in [0.35, 0.92]$
        \item Quantifiers (all → rarely): $B \in [0.22, 0.98]$
        \item Hedging expressions: $\Delta B \in [-0.25, -0.05]$
        \item Intensifiers/Downtoners: $\Delta B \in [-0.15, +0.18]$
    \end{itemize}

    \item \textbf{Composite B-Calculation}: Additive model with floor (0) and ceiling (1) bounds.

    \item \textbf{Discipline Variation}: B-values require discipline-specific calibration.

    \item \textbf{Validation Status}: Ordinal rankings are empirically grounded (T1--T2); cardinal values are proposals (T4--T5) requiring Monte-Carlo validation.

    \item \textbf{Practical Guidelines}: Step-by-step assessment procedure with common pitfalls.
\end{enumerate}

\begin{eslbox}
\textbf{Chapter 4 Self-Assessment}

This chapter presents linguistic marker categories ($B \approx 0.75$) based on established corpus linguistics ($E \approx 0.80$ for existence of markers) but with proposed cardinal values ($E \approx 0.35$ for specific numbers).

\textbf{Weighted Assessment}:
\begin{itemize}
    \item Structure and categories: $K = 0.94$ (T2)
    \item Specific B-values: $K = 0.47$ (T5)
    \item Weighted average: $K \approx 0.70$
\end{itemize}

\textbf{Status}: \textcolor{eslorange}{\textbf{Partially Stable}} ($K \in [0.50, 0.75)$)

\textbf{Honest Assessment}: This chapter's specific numerical values are among the least validated components of ESL. The framework's utility does not depend on exact B-values---ordinal correctness (``proves'' $>$ ``suggests'' $>$ ``hints'') is sufficient for most applications. Cardinal precision requires the corpus validation studies outlined in Section \ref{sec:validation}.
\end{eslbox}

% =============================================================================
% REFERENCES
% =============================================================================

\newpage
\bibliographystyle{apalike}

\begin{thebibliography}{99}

\bibitem[Biber et al., 1999]{Biber1999}
Biber, D., Johansson, S., Leech, G., Conrad, S., \& Finegan, E. (1999).
\textit{Longman Grammar of Spoken and Written English}.
Pearson Education.

\bibitem[Coates, 1983]{Coates1983}
Coates, J. (1983).
\textit{The Semantics of the Modal Auxiliaries}.
Croom Helm.

\bibitem[Hyland, 1998]{Hyland1998}
Hyland, K. (1998).
\textit{Hedging in Scientific Research Articles}.
John Benjamins.

\bibitem[Palmer, 2001]{Palmer2001}
Palmer, F. R. (2001).
\textit{Mood and Modality} (2nd ed.).
Cambridge University Press.

\bibitem[Salager-Meyer, 1994]{SalagerMeyer1994}
Salager-Meyer, F. (1994).
Hedges and textual communicative function in medical English written discourse.
\textit{English for Specific Purposes}, 13(2), 149--170.

\bibitem[Varttala, 2001]{Varttala2001}
Varttala, T. (2001).
\textit{Hedging in Scientifically Oriented Discourse: Exploring Variation According to Discipline and Intended Audience}.
University of Tampere.

\end{thebibliography}

\end{document}
